\section*{Problem Set 4 due on Due March 18 at 11:59pm}
\question{1}{}
Consider a 2D space with negative curvature. Show that the embedding into Euclidean space doe not yield a space with constant curvature for either hyperpoloid:
\begin{itemize}
    \item [(a)] $x^2+y^2-z^2=R^2$
    \item [(b)] $x^2+y^2-z^2=-R^2$
\end{itemize}
Find the Ricci scalar for both cases. Then show the case (b) embedded in a non-Euclidean space does give a space of constant curvature and find the Ricci scalar. Also, what happens in case (a) for a non-Euclidean embedding?
\answer{}
\begin{itemize}
    \item [(a)]
\end{itemize}
For the hyperboloid $x^2+y^2-z^2=R^2$, we can parametrize the surface using cylindrical coordinates ($z, \phi$):
\begin{align}
    x=& \sqrt{R^2+z^2} \cos\phi,\\
    y=& \sqrt{R^2+z^2} \sin\phi,\\
    z=&z.
\end{align}
The original metric in the embedding space is given by
\begin{align}
    ds^2=dx^2+dy^2+dz^2.
\end{align}
Then we have 
\begin{align}
    dx=& \frac{z}{\sqrt{R^2+z^2}} \cos\phi dz - \sqrt{R^2+z^2} \sin\phi d\phi,\\
    dy=& \frac{z}{\sqrt{R^2+z^2}}\sin\phi dz + \sqrt{R^2+z^2} \cos\phi d\phi,\\
    dz=&dz.
\end{align}
Now we can compute the metric on the surface by substituting these expressions into the original metric:
\begin{align}
    ds^2=&dx^2+dy^2+dz^2\\
    =&\left(\frac{z^2}{R^2+z^2}+1\right)dz^2+(R^2+z^2)d\phi^2\\
    =&\frac{R^2+2z^2}{R^2+z^2}dz^2+(R^2+z^2)d\phi^2.
\end{align}
Then, we can say in the coordinates $(z, \phi)$, the metric is given by
\begin{align}
    g_{zz}=&\frac{R^2+2z^2}{R^2+z^2},\\
    g_{\phi\phi}=&R^2+z^2,\\
    g_{z\phi}=&g_{\phi z}=0.
\end{align}
After calculation in Mathematica (see the attached file), we can find the Ricci scalar is given by
\begin{align}
    \mathcal{R}=-\frac{2R^2}{(R^2+2z^2)^2}.
\end{align}
\begin{itemize}
    \item [(b)]
\end{itemize}
For the hyperboloid $x^2+y^2-z^2=-R^2$, we can parametrize the surface using cylindrical coordinates ($z, \phi$):
\begin{align}
    x=& \sqrt{z^2-R^2} \cos\phi,\\
    y=& \sqrt{z^2-R^2} \sin\phi,\\
    z=&z.
\end{align} 
The original metric in the embedding space is given by
\begin{align}
    ds^2=dx^2+dy^2+dz^2.
\end{align} Then we have 
\begin{align}   
    dx=& \frac{z}{\sqrt{z^2-R^2}} \cos\phi dz - \sqrt{z^2-R^2} \sin\phi d\phi,\\
    dy=& \frac{z}{\sqrt{z^2-R^2}}\sin\phi dz + \sqrt{z^2-R^2} \cos\phi d\phi,\\
    dz=&dz.
\end{align}
Now we can compute the metric on the surface by substituting these expressions into the original metric:
\begin{align}
    ds^2=&dx^2+dy^2+dz^2\\
    =&\left(\frac{z^2}{z^2-R^2}+1\right)dz^2+(z^2-R^2)d\phi^2\\ 
    =&\frac{2z^2-R^2}{z^2-R^2}dz^2+(z^2-R^2)d\phi^2.
\end{align}
Then, we can say in the coordinates $(z, \phi)$, the metric is given by
\begin{align}
    g_{zz}=&\frac{2z^2-R^2}{z^2-R^2},\\
    g_{\phi\phi}=&z^2-R^2,\\
    g_{z\phi}=&g_{\phi z}=0.
\end{align}
After calculation in Mathematica (see the attached file), we can find the Ricci scalar is given by
\begin{align}
    \mathcal{R}=\frac{2R^2}{(2z^2-R^2)^2}.
\end{align}
Now, we can see that the Ricci scalar is not a constant, which means the curvature is not constant. However, if we embed the hyperboloid into a non-Euclidean space with metric
\begin{align}
    ds^2=dx^2+dy^2-dz^2,
\end{align}we can find the metric on the surface is given by
\begin{align}
    ds^2=&dx^2+dy^2-dz^2\\
    =&\left(\frac{z^2}{z^2-R^2}-1\right)dz^2+(z^2-R^2)d\phi^2\\
    =&\frac{R^2}{z^2-R^2}dz^2+(z^2-R^2)d\phi^2.
\end{align}
Then, we can say in the coordinates $(z, \phi)$, the metric is given by
\begin{align}
    g_{zz}=&\frac{R^2}{z^2-R^2},\\
    g_{\phi\phi}=&z^2-R^2,\\
    g_{z\phi}=&g_{\phi z}=0.
\end{align}
After calculation in Mathematica (see the attached file), we can find the Ricci scalar is given by
\begin{align}
    \mathcal{R}=-\frac{2}{R^2}.
\end{align}
We can see that the Ricci scalar is a constant, which means the curvature is constant. 

\begin{itemize}
    \item [(a)]
\end{itemize}
For the case (a), if we embed the hyperboloid into a non-Euclidean space with metric
\begin{align}
    ds^2=dx^2+dy^2-dz^2,
\end{align}we can find the metric on the surface is given by
\begin{align}
    ds^2=&dx^2+dy^2-dz^2\\
    =&\left(\frac{z^2}{R^2+z^2}-1\right)dz^2+(R^2+z^2)d\phi^2\\
    =&-\frac{R^2}{R^2+z^2}dz^2+(R^2+z^2)d\phi^2.
\end{align}
Then, we can say in the coordinates $(z, \phi)$, the metric is given by
\begin{align}
    g_{zz}=&-\frac{R^2}{R^2+z^2},\\
    g_{\phi\phi}=&R^2+z^2,\\
    g_{z\phi}=&g_{\phi z}=0.
\end{align}
After calculation in Mathematica (see the attached file), we can find the Ricci scalar is given by
\begin{align}
    \mathcal{R}=\frac{2}{R^2}.
\end{align}
We can see that the Ricci scalar is a constant, which means the curvature is constant. 

\qed




\clearpage
\question{2}{}
Consider the metric
\begin{align}
    g_{\mu\nu}=C_{\mu\nu} + \frac{K}{1-KC_{ab} x^a x^b} C_{\mu\lambda} C_{\nu\sigma} x^\lambda x^\sigma
\end{align}
\begin{itemize}
    \item [(a)] Derive the inverse metric $g^{\mu\nu}$.
    \item [(b)] Derive the Christoffel symbols $\Gamma^a_{bc}$.
    \item [(c)] Derive solutions for the equations of motions when $K=+1,0,-1$.
\end{itemize}
\answer{}
\begin{itemize}
    \item [(a)]
\end{itemize}
To derive the inverse metric $g^{\mu\nu}$, we start with the given metric:
\begin{align}
    g_{\mu\nu}=C_{\mu\nu} + \frac{K}{1-KC_{ab} x^a x^b} C_{\mu\lambda} C_{\nu\sigma} x^\lambda x^\sigma.
\end{align}
We can rewrite this as:
\begin{align}
    g_{\mu\nu} = C_{\mu\nu} + \frac{K}{1-KC_{ab} x^a x^b} (C_{\mu\lambda} x^\lambda)(C_{\nu\sigma} x^\sigma).
\end{align}
Let us define $x_\mu = C_{\mu\lambda} x^\lambda$. Also, we can define $C_{ab} x^a x^b = x^\rho x_\rho$ since $x^\rho = C^{\rho\lambda} x_\lambda$. Then we can express the metric as:
\begin{align}
    g_{\mu\nu} = C_{\mu\nu} + \frac{K}{1-Kx^\rho x_\rho} x_\mu x_\nu.
\end{align}
To find the inverse metric $g^{\mu\nu}$, we can assume it has the form:
\begin{align}
    g^{\mu\nu} = C^{\mu\nu} + \alpha x^\mu x^\nu,
\end{align}
where $\alpha$ is a constant to be determined. We need to satisfy the condition:
\begin{align}
    g^{\mu\rho} g_{\rho\nu} = \delta^\mu{}_\nu.
\end{align}
Substituting the expressions for $g_{\mu\nu}$ and $g^{\nu\rho}$, we get:
\begin{align}
    (C^{\mu\rho} + \alpha x^\mu x^\rho)(C_{\rho\nu} + \frac{K}{1-Kx^\rho x_\rho} x_\rho x_\nu) = \delta^\mu{}_\nu.
\end{align}
Expanding this expression, we have:
\begin{align}
    C^{\mu\rho} C_{\rho\nu} + \frac{K}{1-Kx^\rho x_\rho} C^{\mu\rho} x_\rho x_\nu + \alpha x^\mu x^\rho C_{\rho\nu} + \alpha \frac{K}{1-Kx^\rho x_\rho} x^\mu x^\rho x_\rho x_\nu = \delta^\mu{}_\nu.
\end{align}
Since $C^{\mu\rho} C_{\rho\nu} = \delta^\mu{}_\nu$, we can simplify this to:
\begin{align}
     0&=\frac{K}{1-Kx^\rho x_\rho} C^{\mu\rho} x_\rho x_\nu + \alpha x^\mu x_\nu + \alpha \frac{K}{1-Kx^\rho x_\rho} x^\mu x^\rho x_\rho x_\nu \\
     &=\frac{K}{1-Kx^\rho x_\rho} x^\mu x_\nu + \alpha x^\mu x_\nu + \alpha \frac{K}{1-Kx^\rho x_\rho} x^\mu (x^\rho x_\rho) x_\nu \\
    &=\left(\frac{K}{1-Kx^\rho x_\rho} + \alpha + \alpha \frac{K}{1-Kx^\rho x_\rho} (x^\rho x_\rho)\right) x^\mu x_\nu.
\end{align}
Hence, we can solve for $\alpha$:
\begin{align}
    &\alpha \left(1 + \frac{K x^\rho x_\rho}{1-Kx^\rho x_\rho}\right) = -\frac{K}{1-Kx^\rho x_\rho}\\
    \implies& \alpha \frac{1}{1-Kx^\rho x_\rho} = -\frac{K}{1-Kx^\rho x_\rho}\\
    \implies &\alpha = -K.
\end{align}
Thus, the inverse metric is given by:
\begin{align}
    g^{\mu\nu} = C^{\mu\nu} - K x^\mu x^\nu.
\end{align}
\begin{itemize}
    \item [(b)]
\end{itemize}
First, we need to compute the derivatives of the metric components. The metric is given by:
\begin{align}
    g_{\mu\nu} = C_{\mu\nu} + \frac{K}{1-Kx^\rho x_\rho} x_\mu x_\nu= C_{\mu\nu} + \frac{K}{1-KC_{ab} x^a x^b} C_{\mu\lambda} C_{\nu\sigma} x^\lambda x^\sigma=C_{\mu\nu} +f(x) x_\mu x_\nu,
\end{align}
where we have defined $f(x) = \frac{K}{1-Kx^\rho x_\rho}$. Then, we can compute the derivatives of the metric components:
\begin{align}
    \partial_\lambda g_{\mu\nu} =& \underbrace{\partial_\lambda C_{\mu\nu}}_{=0} + \partial_\lambda f(x) x_\mu x_\nu + f(x) \partial_\lambda (x_\mu x_\nu)\\
    =& \partial_\lambda f(x) x_\mu x_\nu + f(x) (\partial_\lambda x_\mu x_\nu + x_\mu \partial_\lambda x_\nu)\\
    =& \partial_\lambda f(x) x_\mu x_\nu + f(x) (\delta_\lambda{}_\mu x_\nu + x_\mu \delta_\lambda{}_\nu).
\end{align}
Also, the $\partial_\lambda f(x)$ can be computed as:
\begin{align}
    \partial_\lambda f(x) =& \partial_\lambda \left(\frac{K}{1-Kx^\rho x_\rho}\right)\\
    =& \frac{K^2 \partial_\lambda (x^\rho x_\rho)}{(1-Kx^\rho x_\rho)^2}\\
    =& \frac{K^2 \partial_\lambda (C^{\rho\sigma} x_\sigma x_\rho)}{(1-Kx^\rho x_\rho)^2}\\
    =& \frac{K^2 (C^{\rho\sigma} \delta_\lambda{}_\sigma x_\rho + C^{\rho\sigma} x_\sigma \delta_\lambda{}_\rho)}{(1-Kx^\rho x_\rho)^2}\\
    =& \frac{2K^2 x_\lambda}{(1-Kx^\rho x_\rho)^2}.
\end{align}
Hence the derivatives of the metric components can be expressed as:
\begin{align}
    \partial_\lambda g_{\mu\nu} =& \frac{2K^2 x_\lambda}{(1-Kx^\rho x_\rho)^2} x_\mu x_\nu + f(x) (\delta_\lambda{}_\mu x_\nu + \delta_\lambda{}_\nu x_\mu ).
\end{align}

Now we can compute the Christoffel symbols using the formula:
\begin{align}
    \Gamma^a_{bc} = \frac{1}{2} g^{ad} (\partial_b g_{cd} + \partial_c g_{bd} - \partial_d g_{bc}).
\end{align}
Substituting the expressions for $g^{ad}$ and $\partial_\lambda g_{\mu\nu}$, we can compute the Christoffel symbols: 
\begin{align}
    \Gamma^a_{bc} =& \frac{1}{2} (C^{ad} - K x^a x^d) \Bigg( \frac{2K^2 x_b}{(1-Kx^\rho x_\rho)^2} x_c x_d + f(x) (\delta_b{}_{c} x_d + \textcolor{red}{\cancel{\delta_b{}_{d} x_c}} ) \\
    +& \cancel{\frac{2K^2 x_c}{(1-Kx^\rho x_\rho)^2} x_b x_d} + f(x) (\delta_c{}_{b} x_d + \textcolor{blue}{\cancel{\delta_c{}_{d} x_b}} ) - \cancel{\frac{2K^2 x_d}{(1-Kx^\rho x_\rho)^2} x_b x_c} - f(x) (\textcolor{red}{\cancel{\delta_d{}_{b} x_c}} + \textcolor{blue}{\cancel{\delta_d{}_{c} x_b}} )\Bigg)\\
    =& \frac{1}{2} (C^{ad} - K x^a x^d) \Bigg( \frac{2K^2 }{(1-Kx^\rho x_\rho)^2}x_b x_c x_d + 2\frac{K}{1-Kx^\rho x_\rho} \delta_b{}_{c} x_d \Bigg)\\
    =& \frac{K}{1-Kx^\rho x_\rho} (C^{ad} - K x^a x^d) \left( \frac{K}{1-Kx^\rho x_\rho} x_b x_c x_d + \delta_b{}_{c} x_d \right)\\    
    =& \frac{K}{1-Kx^\rho x_\rho} \left( \textcolor{red}{\frac{K}{1-Kx^\rho x_\rho} x_b x_c x^a} + \delta_b{}_{c} x^a - \textcolor{red}{\frac{K^2}{1-Kx^\rho x_\rho} x^a x^d x_d x_b x_c} - K\delta_b{}_{c} x^a x^d x_d \right)\\
    =& \frac{K}{1-Kx^\rho x_\rho} x^a \Bigg( \textcolor{red}{\frac{K}{1-Kx^\rho x_\rho}x_b x_c\left( 1 -Kx^d x_d \right)} + \delta_b{}_{c}(1-Kx^d x_d) \Bigg)\\
    =& Kx^a \Bigg( \frac{K}{1-Kx^\rho x_\rho}x_b x_c + \delta_{bc}  \Bigg).
\end{align}
Note that 
\begin{align}
    \delta_{bc}=\partial_b x_c = =\partial_b (C_{c\lambda} x^\lambda) = C_{c\lambda} \delta_b{}^\lambda = C_{cb}.
\end{align}
Then we can express the Christoffel symbols as:
\begin{align}
    \Gamma^a_{bc} = Kx^a \Bigg( \frac{K}{1-Kx^\rho x_\rho}x_b x_c + C_{bc}  \Bigg)=K x^a g_{bc}.
\end{align}

\begin{itemize}
    \item [(c)]
\end{itemize}
The equations of motion for a geodesic are given by:
\begin{align}
    \frac{d^2 x^a}{d\lambda^2} + \Gamma^a_{bc} \frac{dx^b}{d\lambda} \frac{dx^c}{d\lambda} = 0,
\end{align}
where $\lambda$ is the affine parameter. Substituting the expression for $\Gamma^a_{bc}$, we have:
\begin{align}
    \frac{d^2 x^a}{d\lambda^2} + K x^a g_{bc} \frac{dx^b}{d\lambda} \frac{dx^c}{d\lambda} = 0.
\end{align}
This is a second-order ordinary differential equation for $x^a(\lambda)$. The solutions will depend on the value of $K$:
\begin{itemize}
    \item For $K=0$, the equation reduces to $\frac{d^2 x^a}{d\lambda^2} = 0$, which implies that $v^a= \frac{dx^a}{d\lambda}$ is constant. Thus, the solution is given by:
    \begin{align}
        x^a(\lambda) = v^a \lambda + x_0^a,
    \end{align}
    where $x_0^a$ is the initial position.
    \item For $K=+1$, the equation becomes $\frac{d^2 x^a}{d\lambda^2} + x^a g_{bc} \frac{dx^b}{d\lambda} \frac{dx^c}{d\lambda} = 0$. We can apply $\frac{d}{d\lambda}(g_{ab} v^a v^b)$ to find that $g_{ab} v^a v^b$ is a constant. Let $g_{ab} v^a v^b = \beta$, then we can express the equation as $\frac{d^2 x^a}{d\lambda^2} + \beta x^a = 0$. The solution to this equation is given by:
    \begin{align}
        x^a(\lambda) = A^a \cos(\sqrt{\beta} \lambda) + B^a \sin(\sqrt{\beta} \lambda),
    \end{align}
    where $A^a$ and $B^a$ are constants determined by the initial conditions.
    \item For $K=-1$, the equation becomes $\frac{d^2 x^a}{d\lambda^2} - x^a g_{bc} \frac{dx^b}{d\lambda} \frac{dx^c}{d\lambda} = 0$. Similar to the case of $K=+1$, we can find that $g_{ab} v^a v^b$ is a constant. Let $g_{ab} v^a v^b = \beta$, then we can express the equation as $\frac{d^2 x^a}{d\lambda^2} - \beta x^a = 0$. The solution to this equation is given by:
    \begin{align}
        x^a(\lambda) = C^a e^{\sqrt{\beta} \lambda} + D^a e^{-\sqrt{\beta} \lambda},
    \end{align}
    where $C^a$ and $D^a$ are constants determined by the initial conditions.
\end{itemize}
\qed

\clearpage
\question{3}{}
Problem~8.2 in Carroll. Note that in the 1st edition, there were two errors. As in class, the metric should contain $e^{2Ht}$ instead of $e^{Ht}$. The problem should be solved for non comoving observers. That is, you can not take $dx^i/d\lambda=0$.

Consider the de Sitter space in coordinates where the metric takes the form 
\begin{align}
    ds^2=dt^2-e^{2Ht}(dx^2+dy^2+dz^2)
\end{align} 
Solve the geodesic equation to find the affine parameter as a function of $t$. Show that geodesics reach $t\to-\infty$ in a finite affine parameter, demonstrating that these coordinates fail to cover the entire manifold.
\answer{}
The Christoffel symbols for the given metric can be computed as:
\begin{align}
    \Gamma^\lambda_{\mu\nu}=\frac{1}{2} g^{\lambda\sigma} (\partial_\mu g_{\nu\sigma} + \partial_\nu g_{\mu\sigma} - \partial_\sigma g_{\mu\nu}).
\end{align}
Since the metric is diagonal, we can compute the non-zero Christoffel symbols as:
\begin{align}
    \Gamma^t_{ij}=&\frac{1}{2} g^{tt} (\partial_i g_{jt} + \partial_j g_{it} - \partial_t g_{ij})\\
    =&\frac{1}{2} (1) (0 + 0 - \partial_t (-e^{2Ht} \delta_{ij}))\\
    =&H e^{2Ht} \delta_{ij},
\end{align}
and
\begin{align}
    \Gamma^i_{jt}=&\frac{1}{2} g^{ik} (\partial_j g_{kt} + \partial_t g_{jk} - \partial_k g_{jt})\\
    =&\frac{1}{2} (-e^{-2Ht} \delta^{ik}) (0 + \partial_t (-e^{2Ht} \delta_{jk}) - 0)\\
    =&H \delta^i_j.
\end{align}
The geodesic equations are given by:
\begin{align}
    \frac{d^2 x^\lambda}{d\lambda^2} + \Gamma^\lambda_{\mu\nu} \frac{dx^\mu}{d\lambda} \frac{dx^\nu}{d\lambda} = 0.
\end{align}
Substituting the expressions for the Christoffel symbols, we can write the geodesic equations as:
\begin{align}
    \frac{d^2 t}{d\lambda^2} + H e^{2Ht} \delta_{ij} \frac{dx^i}{d\lambda} \frac{dx^j}{d\lambda} = 0,
\end{align}and
\begin{align}
    \frac{d^2 x^i}{d\lambda^2} + 2H \frac{dt}{d\lambda} \frac{dx^i}{d\lambda} = 0.
\end{align}
We can rewrite the second equation as:
\begin{align}
    \frac{d}{d\lambda} \left( e^{2Ht} \frac{dx^i}{d\lambda} \right) = 0\\
    \implies e^{2Ht} \frac{dx^i}{d\lambda} = C^i,
\end{align}
where $C^i$ are constants determined by the initial conditions. Substituting this expression into the first equation, we have:
\begin{align}
    \frac{d^2 t}{d\lambda^2} + H e^{2Ht} \delta_{ij} \frac{C^i}{e^{2Ht}} \frac{C^j}{e^{2Ht}} = 0\\
    \implies \frac{d^2 t}{d\lambda^2} + H e^{-2Ht} \delta_{ij} C^i C^j = 0\\
    \implies \frac{d^2 t}{d\lambda^2} + H e^{-2Ht} C^2 = 0,
\end{align}
where $C^2 = \delta_{ij} C^i C^j$. This is a second-order ordinary differential equation for $t(\lambda)$. We consider the following equation:
\begin{align}
    &\frac{d}{d\lambda} \left( \Big( \frac{d t}{d\lambda} \Big)^2 -  e^{-2Ht} C^2\right)\\
    =& 2 \frac{d t}{d\lambda} \frac{d^2 t}{d\lambda^2} + 2H e^{-2Ht} C^2 \frac{dt}{d\lambda}\\
    =& 2 \frac{d t}{d\lambda} \left( \frac{d^2 t}{d\lambda^2} + H e^{-2Ht} C^2 \right)\\
    =& 0.
\end{align}
Hence, we can find that $\Big( \frac{d t}{d\lambda} \Big)^2 -  e^{-2Ht} C^2$ is a constant. Let $\Big( \frac{d t}{d\lambda} \Big)^2 -  e^{-2Ht} C^2 = \beta$, then we can express the equation as:
\begin{align}
    &\Big( \frac{d t}{d\lambda} \Big)^2 = \beta + e^{-2Ht} C^2\\
    \implies &\frac{d t}{d\lambda} = \sqrt{\beta + e^{-2Ht} C^2}.
\end{align}

\textbf{If $\beta=0$}, we can express the equation as:
\begin{align}
    &\frac{d t}{d\lambda} = \sqrt{e^{-2Ht} C^2}\\
    \implies &\frac{d t}{d\lambda} = e^{-Ht} C\\
    \implies &\lambda = -\frac{1}{H} e^{Ht} + \lambda_0,
\end{align}
where $\lambda_0$ is a constant determined by the initial conditions. We can see that as $t \to -\infty$, $\lambda \to \lambda_0$, which means that geodesics reach $t \to -\infty$ in a finite affine parameter. This demonstrates that these coordinates fail to cover the entire manifold.

\textbf{If $\beta \neq 0$}, I don't know how to solve the equation analytically, but we can still find that as $t \to -\infty$, $\lambda$ approaches a finite value, which also demonstrates that these coordinates fail to cover the entire manifold.
\qed